\newpage
\section{问题重述}

\subsection{问题背景}
城市交通拥堵严重影响居民出行效率和城市经济运行。智能交通系统通过实时感知交通状态、预测短时流量、优化信号控制，可以有效缓解交通拥堵。本文研究基于图神经网络的短时交通流预测与信号协同控制问题。

\subsection{问题提出}
\textbf{问题一}：构建城市交通网络的图结构表示，分析交通流的时空相关性特征。

\textbf{问题二}：基于GraphSAGE等图神经网络模型，建立短时交通流预测模型（预测未来15分钟流量）。

\textbf{问题三}：设计多路口协同信号控制策略，以总延误时间最小为目标，考虑绿波带约束和相位差优化。

\subsection{研究意义}
为城市交通精细化管理和信号配时优化提供数据驱动的方法论支撑。
